% RQ5 Abstract: Inter-DAO Cooperation

As the DAO ecosystem matures, organizations increasingly face challenges that transcend individual communities. Public goods funding, protocol interoperability, and ecosystem-wide standards require coordination across organizational boundaries. Yet inter-DAO cooperation mechanisms remain underdeveloped.

We investigate how cross-DAO collaboration affects governance outcomes using multi-agent simulation. Through \experimentruns{} simulation runs across five cooperation scenarios, we examine joint proposal mechanisms, resource sharing arrangements, and cross-DAO voting coordination.

Our results reveal that inter-DAO cooperation can generate positive-sum outcomes when participating DAOs have complementary resource profiles. Homogeneous DAOs (similar size, preferences) achieve higher cooperation success rates but smaller gains. Heterogeneous configurations unlock larger benefits but face coordination challenges. Asymmetric relationships (large-small DAO pairings) require explicit fairness mechanisms to remain stable.

We identify conditions under which cooperation emerges naturally versus requiring explicit incentive design, and provide guidelines for inter-DAO proposal structures. Our findings inform the emerging practice of DAO-to-DAO relations and multi-organization governance.
